% =====================================================================
% Section 4: Method — EvidenceFlow
% =====================================================================
%
% Status: Draft §4 Method v3 — opening overview + 3 mechanism subsections (2026-05-15)
%   v1: 5 subsections — backed up at archive/04_method_v1_5sections_20260515.tex
%   v2: 3 subsections — collapsed §4.3+§4.4+§4.5 into "Query-Local Evidence Flow"
%       to remove pipeline visual cue while keeping all paragraph-level content
%   v3: unnumbered method overview paragraph + three numbered mechanism subsections
%       following the Methodology structure of PropRAG / Relink / NeocorRAG
% Framing reference: research_memory/emnlp_expand_then_compose/46_*.md
% Writing discipline: research_memory/emnlp_expand_then_compose/45_*.md
%
% Method name: \methodname (defined via \newcommand in main file)
%
% Discipline checklist:
%   - No internal codename anywhere in main paper
%   - No "stage / pool / module / readout / selector / composer"
%     as method-internal terms (only allowed in reverse-framing sentences)
%   - No "objective / certify / unified" — use
%     "criterion / boundary / proxy / formation principle"
%   - 4.x chapter titles must be mechanism names, not module names
%
% =====================================================================

\section{Method: \methodname{}}
\label{sec:method}

In this section, we introduce \methodname{}, which treats multi-hop retrieval
as the construction of a document-level evidence chain rather than the
independent ranking of documents. Given a question $q$, a corpus
$\mathcal{D}$, a reader budget $K$, and a pre-extracted set of OpenIE facts
$\mathcal{F}$, \methodname{} returns an evidence set
$E_K \subseteq \mathcal{D}$ by constructing a query-local document graph
$G_q=(V_q,A_q)$ and terminating it at the reader boundary. First, we describe
the offline document-grounded OpenIE graph that supplies entity anchors and
relational structure (\S\ref{sec:method:substrate}). Second, we show how
\methodname{} forms a
query-local evidence graph whose vertices and arcs are activated under $q$
(\S\ref{sec:method:flow}). Finally, we define the chain-integrity boundary
that forms $E_K$ under budget (\S\ref{sec:method:boundary}). Chain
connectivity is treated throughout as a property of the retrieval graph
itself, not as a relation reconstructed downstream over a flat passage pool.

\subsection{Offline Indexing: Document-Grounded OpenIE Graph}
\label{sec:method:substrate}

We build a document-grounded graph
$\mathcal{G}_{\mathcal{D}}=(\mathcal{V}_{\mathcal{D}},\mathcal{E}_{\mathcal{D}})$
over the corpus $\mathcal{D}$ during offline indexing. The graph supplies
the entity anchors and relational structure that \methodname{} draws on at
query time (\S\ref{sec:method:flow}). Following standard OpenIE practice
\citep{angeli2015leveraging,gutierrez2025hipporag}, we extract salient
assertions from each document and organize them into the graph below; the
extractor itself is not a contribution of this work.

\paragraph{OpenIE fact extraction.}
For each document $d \in \mathcal{D}$, a language model extracts a set of
OpenIE facts $f = (h, r, t)$, where $h$ and $t$ are head and tail entities
and $r$ is the relation phrase that connects them. We retain each fact
together with its provenance $\mathrm{src}(f) \in \mathcal{D}$, so that
every fact can be traced back to the document from which it was extracted.
The resulting fact set is denoted $\mathcal{F}$.

\paragraph{Graph construction.}
The graph $\mathcal{G}_{\mathcal{D}}$ contains two node types and three
edge types. Nodes are: (i) \emph{passage nodes}, one per document
$d \in \mathcal{D}$, carrying the natural-language content that a
downstream reader will eventually consume; and (ii) \emph{phrase nodes},
one per distinct entity surface form extracted across $\mathcal{F}$.
Edges are: (i) \emph{passage--phrase} edges that connect each document to
the entities it mentions; (ii) \emph{relational} edges that connect a pair
of phrase nodes $(h, t)$ whenever some fact $f = (h, r, t) \in \mathcal{F}$
asserts a relation between them, with $r$ retained as an edge attribute
and $\mathrm{src}(f)$ retained as the edge provenance; and (iii)
\emph{synonymy} edges that connect phrase nodes whose dense embeddings
are highly similar, providing alias coverage across surface variants of
the same entity.

\paragraph{Indexing for retrieval.}
We pre-compute dense embeddings for all passage and phrase nodes, and
materialize $\mathcal{G}_{\mathcal{D}}$ as an adjacency index keyed by
node identifier. The graph is built once and reused across all questions:
at query time, no new OpenIE extraction or fact generation is performed
over $\mathcal{D}$. The artifacts produced by this stage---phrase nodes,
passage--phrase links, relational edges with provenance, and synonymy
edges---are exactly the pieces \methodname{} draws on when it constructs
a query-local subgraph and ranks documents by graph-grounded propagation
(\S\ref{sec:method:flow}).

\subsection{Query-Local Evidence Graph}
\label{sec:method:flow}

Given the offline graph $\mathcal{G}_{\mathcal{D}}$, \methodname{} constructs the query-local
graph through one continuous evidence-flow process. The process activates a
compact region of documents from question anchors and entry signals, then
admits document-to-document transitions when they are witnessed by facts in
$\mathcal{F}$. We describe these operations as facets of the same graph
formation process rather than as separate retrieval and selection stages.

\paragraph{Question-anchored activation.}
\label{sec:method:anchor}
Evidence flow begins by activating a compact document region
$V_q \subseteq \mathcal{D}$ around the question. Activation combines two
complementary signals: a \emph{question anchor} $A(q) \subseteq \mathcal{E}$
extracted from $q$, where $\mathcal{E}$ is the entity vocabulary induced by
$\mathcal{F}$, and an \emph{entry signal} $\mathrm{Entry}(q) \subseteq
\mathcal{D}$ obtained from a standard dense or lexical retriever. The
anchor identifies which entities $q$ is asking about; the entry signal
identifies which documents are the most relevant starting points in
$\mathcal{D}$. Neither signal alone determines the evidence delivered to
the reader: $V_q$ is initialized from these two signals and then formed
where they can be connected through facts in $\mathcal{F}$. Concretely,
$V_q$ is initialized with the union of (i) documents in $\mathrm{Entry}(q)$
and (ii) documents that are the provenance of at least one fact in
$\mathcal{F}$ whose arguments include an anchor entity in $A(q)$. From this
initial vertex set, $V_q$ is expanded through a bounded number of hops
along facts whose entities are connected to those already activated, with
the size of $V_q$ bounded to keep the activated region compact relative to
$\mathcal{D}$, while the reader budget $K$ ultimately constrains $E_K$.
This activation rule keeps $|V_q|$ small enough to make downstream
computation tractable, while still allowing documents that are reachable
only through multi-hop entity chains to enter $V_q$ even when they have
low surface similarity to $q$. The entry signal seeds activation but does
not by itself determine $E_K$: documents in $\mathrm{Entry}(q)$ enter
$V_q$ for consideration, yet their position in the entry ranking is not
the sole criterion that decides which of them reach the reader. The anchor
extraction is the mechanism by which $q$ exerts query-conditional control
over $G_q$: anchors are extracted from entity mentions and relation-bearing
expressions in $q$, using the same controlled language model used
throughout our protocol. This extraction is applied to the question, not
to the retrieved corpus, and does not generate evidence chains; it provides
the question-side semantic handles that allow $G_q$ to grow toward
documents that $q$ actually asks about, rather than only toward documents
that happen to score well on dense similarity.

\paragraph{Fact-grounded transitions.}
\label{sec:method:transition}
The activated region becomes a document-level evidence graph only when the
arcs $A_q$ are admitted under fact witnesses. The aim is to make chain
connectivity a structural property of $G_q$, audited at the moment of
admission, so that the evidence chain reaching the reader is constructed
within the retrieval graph itself rather than reassembled from text
afterward.

Let $f = (h, r, t) \in \mathcal{F}$ denote a fact extracted from some
source document $\mathrm{src}(f) \in \mathcal{D}$, with head and tail
entities $h, t \in \mathcal{E}$. For an ordered pair of documents
$(d_i, d_j)$, the \emph{witness set} under question $q$ is
\begin{equation}
\label{eq:witness-set}
W_q(d_i, d_j) = \bigl\{\, w \in \mathcal{W}(\mathcal{F}) \;\bigm|\;
w \text{ links } d_i \text{ and } d_j \text{ under } q \,\bigr\},
\end{equation}
where $\mathcal{W}(\mathcal{F})$ ranges over witness objects induced by
facts in $\mathcal{F}$: a single fact $f$ whose endpoints reach across
$d_i$ and $d_j$, or a pair of facts $(f_i, f_j)$ with
$\mathrm{src}(f_i) = d_i$ and $\mathrm{src}(f_j) = d_j$ that share an
activated argument. A witness $w$ links $d_i$ and $d_j$ under $q$ when one
of its arguments is activated by $q$ (either through the anchor set
$A(q)$ or through entities already entered the activated region of $V_q$)
and a complementary argument grounds the other document. An arc
$(d_i, d_j)$ is \emph{eligible} for $A_q$ when $W_q(d_i, d_j)$ is
non-empty; surface similarity between $d_i$ and $d_j$ or their individual
scores against $q$ may seed activation but are not by themselves
sufficient to create an arc.

\paragraph{Grounding modes and auditability.}
Every arc in $A_q$ therefore carries an explicit witness in $\mathcal{F}$:
each chain that ultimately reaches the reader can be traced back, edge by
edge, to the source facts and their provenance in $\mathcal{D}$. This is
what distinguishes \methodname{}'s notion of connectivity from a
similarity graph in which two documents are joined whenever they share
topic vectors: similarity may motivate co-occurrence, but it does not
establish that the two documents support a multi-hop step that $q$ asks
about. Witness sets provide this support by construction.

Within Eq.~\eqref{eq:witness-set}, two complementary modes of linkage
naturally arise. A \emph{sentence-grounded witness} is a single fact whose
head and tail are both anchored in a source sentence that already mentions
an entity activated by $q$; this captures cases where one document
explicitly bridges to the other through a $q$-relevant entity within a
bridging sentence. An \emph{argument-grounded witness} is induced by fact
arguments grounded in the two documents through a shared activated
entity, capturing cases where the multi-hop step is realized through
entity coreference across documents rather than a single bridging
sentence. Both modes share the same admission condition (a non-empty
$W_q$) and the same auditability property; they differ only in how the
witness is realized within $\mathcal{F}$.

A crucial property of Eq.~\eqref{eq:witness-set} is that witnesses range
over $\mathcal{W}(\mathcal{F})$, which is induced by the pre-extracted
substrate from \S\ref{sec:method:substrate}, and not over text retrieved
at query time. At the point of arc admission, the query-conditional
information used to test witnesses is grounded in $A(q)$ and the current
state of $V_q$; witnesses themselves are not generated, paraphrased, or
re-extracted from $\mathcal{D}$ for $q$. Chain structure thus arises as a
property of the retrieval graph rather than as a relation reconstructed
from retrieved text at query time.

\subsection{Reader-Boundary Chain Integrity}
\label{sec:method:boundary}
The same evidence-flow process terminates when the reader budget $K$ is
reached: $E_K$ is the budget-bounded outcome of the graph formation, not
an independent ranking step over the activated documents. We formulate
this termination as a \emph{chain-integrity boundary} on $G_q$. $E_K$
should preserve the witnessed connectivity already established in $A_q$,
and any document admitted into $E_K$ should either extend an existing
witnessed chain or address a \emph{question-induced obligation} that no
document already in $E_K$ covers. These obligations --- entity mentions
and relation-bearing slots that $q$ expects the evidence to jointly
resolve --- are extracted from $q$ with the same controlled language model
used for the question anchor; they are applied to the question, not to
the retrieved corpus, and do not generate evidence chains. Concretely,
$E_K$ is formed incrementally from $G_q$ in two interleaved modes.
\emph{Witness-preserving admission} extends $E_K$ along arcs of $A_q$, so
that documents already supported by a witnessed transition to some
$d \in E_K$ are preferred to documents that would only enter $E_K$
through unwitnessed similarity. \emph{Residual admission} is reserved for
cases where an open obligation remains and a candidate document, not yet
connected to $E_K$ by a witnessed transition, provides evidence for it.
Residual admission does not create new unwitnessed arcs: once a document
is admitted, any further movement involving it in $G_q$ remains subject
to the same witness conditions described above.
